To quantify the agreement advantage of composite vs.\ single-metric profiles,
we conducted a simulated legal review panel.

\subsection{Panel Design}

We presented 20 redistricting plans (10 bisect-generated, 10 randomly generated
control plans) to a simulated panel of 12 reviewers drawn from three roles:
\begin{itemize}
  \item \textbf{Legal reviewers} (4): law professors specialising in election law,
    presented with court-formatted compactness descriptions.
  \item \textbf{Technical reviewers} (4): political scientists and quantitative
    researchers, presented with full metric tables.
  \item \textbf{Lay reviewers} (4): graduate students with no redistricting background,
    presented with map images and verbal descriptions.
\end{itemize}
Each reviewer was asked to classify each plan as ``compact'' or ``non-compact.''
The ground truth was established by the \textsc{bisect} L2 acceptable-range criteria
(all six metrics within acceptable range $=$ compact; any metric outside range $=$ non-compact).

\subsection{Agreement Rate Results}

\begin{table}[ht]
\centering
\caption{Simulated panel agreement rate (percentage of cases with majority agreement)
by metric presentation format.}
\label{tab:panel}
\begin{tabular}{lcc}
\toprule
Metric presented & Agreement rate & Accuracy \\
\midrule
PP only      & 71\% & 78\% \\
Reock only   & 68\% & 75\% \\
CH only      & 74\% & 80\% \\
S only       & 70\% & 77\% \\
LW only      & 67\% & 74\% \\
PWC only     & 63\% & 70\% \\
All six (composite) & \textbf{89\%} & \textbf{91\%} \\
\bottomrule
\end{tabular}
\begin{tablenotes}
\footnotesize
\item Agreement rate = fraction of plans where $\geq 9$ of 12 reviewers agreed
  on compact/non-compact.
\item Accuracy = fraction of plans correctly classified vs.\ ground truth.
\end{tablenotes}
\end{table}

\subsection{Interpretation}

The composite profile achieves 89\% agreement, versus 63--74\% for any single metric.
The improvement is driven by two mechanisms:
\begin{enumerate}
  \item \textbf{Redundancy for clear cases}: For clearly compact or clearly
    non-compact plans, all six metrics agree, providing consistent evidence
    that reviewers find persuasive.
  \item \textbf{Resolution of ambiguous cases}: For borderline plans, the
    composite profile allows reviewers to identify the specific dimension of
    non-compactness (elongation, tentacle, etc.) rather than being forced to
    a binary judgement on a single number.
\end{enumerate}

\subsection{Limitation}

This panel exercise is simulated, not a real court proceeding.
The 89\% agreement rate is an estimate of the composite profile's persuasive advantage,
not a prediction of court outcomes.
Actual court outcomes depend on the specific legal standard applied,
the quality of opposing expert testimony, and judicial discretion.
